\input{../tex-inputs/latex-leseansicht-vorspann}

\section[Arthur und Olga Schnitzler an Gustav Schwarzkopf, 1{[}7?{]}. 5. 1912]{L04453 Arthur und Olga Schnitzler an Gustav Schwarzkopf, 1{[}7?{]}. 5. 1912}
\nopagebreak\mylabel{L04453v}
\rehead{ }\normalsize\beginnumbering\briefempfaengerindex{Schwarzkopf, Gustav@\textsc{Schwarzkopf, Gustav}!zzzSchnitzler, Olga@\emph{von Olga Schnitzler}!1912-05-171@{1{[}7?{]}. 5. 1912}|(be}\briefempfaengerindex{Schwarzkopf, Gustav@\textsc{Schwarzkopf, Gustav}!zzzSchnitzler, Arthur@\emph{von Arthur Schnitzler}!1912-05-171@{1{[}7?{]}. 5. 1912}|(be}
\toendnotes[C]{\smallbreak\pagebreak[2]}
\correspDesc{Versand  durch Arthur Schnitzler, Olga Schnitzler am 1[7?]. 5. 1912 in Venedig
\newline{}Übermittlung  am 17. 5. 1912 in Venedig
\newline{}Erhalt  durch Gustav Schwarzkopf im Zeitraum [19. 5. 1912 – 23. 5. 1912?] in Wien}\toendnotes[C]{\smallbreak}
\Standort{DLA, A:Schnitzler, HS.1985.1.1897.}
\physDesc{Bildpostkarte, 92 Zeichen
\newline{}Handschrift Arthur Schnitzler: Bleistift, deutsche Kurrent
\newline{}Handschrift Olga Schnitzler: Bleistift, lateinische Kurrent
\newline{}Versand: Stempel: »\nobreak{}\oindex{Venedig@\textbf{Venedig}|pwk}Venezia, 17. 5. 1912, 22, Ferrovia\nobreak{}«.  }\toendnotes[C]{\smallbreak}\pstart{}{\pb}\textsc{Austria}\oindex{Österreich@\textbf{Österreich}|pw}\pend{}\pstart{}Herrn \textsc{Gustav Schwarzkopf}\pend{}\pstart{}Wien I\oindex{I., Innere Stadt@\textbf{I., Innere Stadt}|pw}\pend{}\pstart{}\textsc{Tiefer Graben 17}\oindex{Wien@\textbf{Wien}!I., Innere Stadt@\textbf{I., Innere Stadt}!Tiefer Graben 17@\textbf{Tiefer Graben 17}|pw}\pend{}{\bigskip}
\pstart
           \noindent{}{\pb}\textcolor{gray}{\textbf{Venezia – Ganal Grande\oindex{Canal Grande@\textbf{Canal Grande}|pw} preso dall’ Accademia\oindex{Gallerie dell’Accademia@\textbf{Gallerie dell’Accademia}|pw} e Palazzo Franchetti\oindex{Palazzo Cavalli-Franchetti@\textbf{Palazzo Cavalli-Franchetti}|pw}}}\pend
           \vspace{1em}
\pstart
           \noindent{}{\pb}Herzliche Grüße!\pend
           
\pstart
           Ihr \spacefill\mbox{Arthur}{\\[\baselineskip]}\spacefill\mbox{{[}hs. O. Schnitzler:{]} Olga.}\pend
           \leftskip=0em{}
\pstart
           {[}hs. A. Schnitzler:{]} \label{K_L04453-1v}\edtext{18. 5. 912}{\lemma{\textnormal{\emph{18. 5. 912}}}\Cendnote{\textnormal{Vermutlich irrte sich Schnitzler im Datum, denn der Poststempel gibt den 17. 5. 1912 an. Denkbar, wenn auch weniger wahrscheinlich, ist daneben die Möglichkeit, dass die Karte tatsächlich am 18. 5. 1912 geschrieben und versandt wurde, der Poststempel jedoch nicht auf den Tag aktualisiert war.}}}\label{K_L04453-1}\pend
           \selectlanguage{ngerman}\endnumbering\briefempfaengerindex{Schwarzkopf, Gustav@\textsc{Schwarzkopf, Gustav}!zzzSchnitzler, Olga@\emph{von Olga Schnitzler}!1912-05-171@{1{[}7?{]}. 5. 1912}|)be}\briefempfaengerindex{Schwarzkopf, Gustav@\textsc{Schwarzkopf, Gustav}!zzzSchnitzler, Arthur@\emph{von Arthur Schnitzler}!1912-05-171@{1{[}7?{]}. 5. 1912}|)be}\mylabel{L04453h}
\begin{anhang}
\end{anhang}\newcommand{\dateiname}{L04453}\newcommand{\titel}{Arthur und Olga Schnitzler an Gustav Schwarzkopf, 1[7?]. 5. 1912}\newcommand{\editorInnen}{Herausgegeben von Jahnke, SelmaMüller, Martin Anton}\input{../tex-inputs/latex-leseansicht-abspann}
